% ---------------------------------------------------------------------------
% Overhead-decomposition analysis for the sync-instrumented runs.
% Drop this into sn-article.tex at the end of Section~\ref{sec:results}
% (just before \section{Related Work}). Requires the figure
% figures/sync_overhead.png (produced by benchmarks/plot_sync.py) and the
% numbers from benchmarks/csv_sync/*.csv (produced by aggregate_sync.py).
% ---------------------------------------------------------------------------

\subsection{Analysis of Overhead}\label{sec:overhead-analysis}

The speedup figures above show \emph{how well} each benchmark scales, but not
\emph{where} the parallel overhead is spent. To decompose it, we re-ran the
three largest workloads --- computational $\pi$ ($10^8$), SumEuler
($4\times10^4$) and matrix multiplication ($4000^2$) --- against an instrumented
copy of the runtime that records, per run, the time each worker spends inside
the user function (\emph{aggregate compute}), the time the coordinator spends
blocked at the \lstinline{syncPList} barrier (\emph{sync wait}), the process
spawn time, the volume of data copied between processes, and the peak
Erlang process-heap footprint sampled every $20$\,ms. Because the instrumentation
itself perturbs timing, we use these runs only for the \emph{ratios} in
\cref{tab:overhead} and \cref{fig:sync-overhead}; the speedups reported earlier come from
the uninstrumented runtime.

\begin{table}
\centering
\resizebox{\columnwidth}{!}{%
\begin{tabular}{|l|c|c|c|c|c|c|}
\hline
Benchmark & Workers & Aggr.\ compute (s) & Sync frac. & Spawn (s) & Msg (MB) & Mem peak (GB) \\
\hline
$\pi$ $10^8$           & 12 & 22.26  & 0.17 & 0.001  & 500.0 & 7.65 \\
SumEuler $4\times10^4$ & 28 & 32.87  & 0.99 & 0.0003 & 0.20  & 0.06 \\
Matmul $4000^2$        & 28 & 8216.3 & 0.88 & 37.36  & 67.8  & 13.46 \\
\hline
\end{tabular}%
}
\caption{Overhead decomposition at each benchmark's peak-speedup worker count
(cf.\ \cref{tab:runtimes}). \emph{Aggr.\ compute} is worker time inside the user
function summed over all workers; \emph{Sync frac.} is coordinator barrier-wait
time as a fraction of wall time; \emph{Spawn} is the total time spent spawning
worker processes per run; \emph{Msg} is data copied between processes per
run; \emph{Mem peak} is the peak process-heap footprint. Means over 10 runs.}
\label{tab:overhead}
\end{table}

\begin{figure}[t]
\centering
\includegraphics[width=0.99\textwidth]{figures/sync_overhead.png}
\caption{Process-memory peak (solid, left axis, note the per-panel scale) and
sync-barrier fraction $\mathit{wait}/\mathit{wall}$ (dashed, right axis) against
worker count for the three instrumented workloads.}
\label{fig:sync-overhead}
\end{figure}

Three observations follow. First, aggregate compute is essentially invariant to
the worker count (e.g.\ SumEuler stays at $27$--$33$\,s across $2$--$28$ workers),
confirming that the \lstinline{PList} model divides a fixed body of work rather
than duplicating it; the mild rise at high worker counts (MatMul grows
$\sim\!19\%$ from $2$ to $28$ workers) is scheduler and garbage-collection
contention. Second, the sync fraction distinguishes the scaling regimes visible
in \cref{fig:sync-overhead}: for SumEuler and MatMul it stays near unity, meaning the
coordinator is a thin barrier and the workers are saturated, whereas for
computational $\pi$ it collapses from $0.61$ at two workers to under $0.09$ by
$28$ --- the coordinator runs out of parallel work to wait for. This, together
with the wall-clock time flattening at $\sim\!10$\,s beyond ten workers, is the
signature of a sequential fraction (range materialisation, the $500$\,MB input
message, and the final reduction) dominating, and accounts for the $\pi$ plateau
noted in \cref{fig:cpi} without appeal to cache behaviour. Third, the copy-based
message-passing model carries a memory cost the speedup plots do not expose:
SumEuler's footprint is negligible ($19$--$63$\,MB) and grows gently with workers,
but MatMul rises to $13.5$\,GB --- each worker closure captures a copy of matrix
$B$, which also drives its spawn time to $37$\,s at $28$ workers --- and
computational $\pi$ sits at $7$--$12$\,GB, dominated by a single $500$\,MB range
message. Reducing these costs, by streaming input rather than materialising and
copying whole \lstinline{Vect}s and by overlapping the final reduction with
computation, is a natural direction for future work.
